\chapter{この研究への含み}
\label{ch:implications}

最後に、ここまでの基本と研究を、LLM-BLM の研究に当てはめる。

\section{LLM-BLM で分かったこと}

\begin{itemize}
\item \SP{} の決算発表文では、LLM（gpt-4o-mini）・FinBERT・ChronoGPT の読み取りは、発表をはさむ 2 日の株価の反応を強く説明した。
\item しかし、その情報は発表の当日にほぼ株価に入りきり、翌日以降のドリフトは有意に予測しなかった。
\item その点数を Black--Litterman の見通しにした月次・日次のロングのみのポートフォリオは、取引コストの後で時価総額加重を上回れなかった。
\item 注目の低い小型株で、文章の読み取りが発表の後 2〜60 取引日のドリフトを予測するかを、事前登録した実験 S1 で確かめている（この報告書を書いた時点で、結果はまだ出ていない）。
\end{itemize}

\section{研究の文献の中での位置}

\begin{table}[H]
\centering\small
\caption{LLM-BLM の結果と、関係する研究}
\label{tab:position}
\begin{tabular}{@{}P{0.32\linewidth}P{0.6\linewidth}@{}}
\toprule
LLM-BLM の結果 & 整合する研究 \\
\midrule
文章の点数は 2 日の反応を強く説明する & Tetlock ほか（2008）、Loughran・McDonald（2011）、Lopez-Lira・Tang（2023）、Chen・Kelly・Xiu（2022）（第 \ref{ch:text} 章） \\
翌日以降のドリフトは残らない & Martineau（2022）：大型株の PEAD は 2006 年以降なく、最近はマイクロキャップでも消えた。Lopez-Lira・Tang：LLM の普及で戦略の収益が低下（第 \ref{ch:event}・\ref{ch:text} 章） \\
時価総額加重に勝てない & ヘッジファンドの平均の手数料の後のアルファは小さい（第 \ref{ch:hfperf} 章）。公表されたシグナルは減衰する（第 \ref{ch:equityfactor}・\ref{ch:failures} 章） \\
小型株ならドリフトが残るかもしれない（S1 の仮説） & Hong・Lim・Stein（2000）、Chordia ほか（2009）、Akbas ほか、PEAD.txt（第 \ref{ch:event} 章）。ただしコストで消えうる（Chordia ほか 2009、Avramov ほか 2023） \\
\bottomrule
\end{tabular}
\end{table}

LLM-BLM の結果は、文献の主な流れとよく一致している。「LLM を使えば市場に勝てる」という期待に対して、「LLM は情報を正しく読むが、その情報はすでに価格に入っている」という、効率的な市場の考え方に沿った答えである。

\section{プロの工程との違い}

\begin{table}[H]
\centering\small
\caption{クオンツの運用の工程と、LLM-BLM の設計の違い}
\label{tab:gap}
\begin{tabular}{@{}P{0.14\linewidth}P{0.27\linewidth}P{0.24\linewidth}P{0.26\linewidth}@{}}
\toprule
工程 & プロの典型 & LLM-BLM & 差の影響 \\
\midrule
シグナル & 数十〜数百を組み合わせ、互いの相関を考える & LLM などの読み取りの一つ & 一つの弱いシグナルでは IR が小さい（基本法則の BR が小さい） \\
速さ & 発表の数ミリ秒〜数分で売買 & 月末・日次の終値 & 情報の大部分は売買の前に価格に入る \\
期間 & シグナルの減衰の曲線に合わせる & 月次・日次・四半期 & 速いシグナルを遅い売買で使うと、価値が消える \\
構築 & ロング・ショート、ファクター中立、TC を高く & ロングのみ、TE 3\% & 悪い読み取りを空売りに使えず、TC が低い \\
コスト & 平方根則のモデル、売買しない領域、執行の最適化 & 片道 0.3\% の一定 & 小さな資金なら妥当。大きな資金では過小 \\
リスク & ファクター型のリスクモデル、ストレス・テスト & Ledoit--Wolf の共分散 & 小規模なら十分 \\
検証 & 長いホールドアウト、試した数の記録、紙の上の運用 & 事前登録（S1）、Holm の補正 & 研究としては厳しい側 \\
\bottomrule
\end{tabular}
\end{table}

\section{研究で試せること}

ここまでの章から、学生の研究の規模で試す価値がありそうなことを、優先度の高い順に挙げる。どれも、新しい分析を始める前に事前登録の追記として書くべきものである。

\begin{enumerate}
\item \textbf{シグナルの価値を、ポートフォリオの制約と分けて測る}。ロングのみ・TE の上限の中での成績だけでなく、紙の上のロング・ショート（上位 − 下位の分位のポートフォリオ）の成績と、その移転係数（TC）を報告する。予測の価値がないのか、ロングのみの制約で価値が失われたのかが分かる（\chapref{ch:perf}、\chapref{ch:pc}）。
\item \textbf{減衰の曲線を描く}。発表の翌日の始値、翌日の終値、1 週間後、1 か月後に売買したとき、それぞれの IC と、コストの後の収益がどう変わるか。情報がどの時間の単位で価格に入るかを示せば、「なぜ月次の運用では勝てないか」を直接に説明できる（\chapref{ch:signal}）。
\item \textbf{既知のファクターへの増分を測る}。LLM の点数を、モメンタム・割安・収益性・決算の驚き（SUE）などに回帰した残差の IC を測る。LLM の点数が、数値の驚きの言い換えにすぎないのか、文章でしか分からない情報を持つのかが分かる（第 \ref{ch:equityfactor}・\ref{ch:event} 章）。
\item \textbf{反応の遅れやすい場面に絞る}。小型株（S1）のほか、アナリストの少ない銘柄、金曜日や発表の多い日、悪い知らせ、長く複雑な文章（ガイダンス、リスクの記述、前回からの文章の変化）。先行研究が遅れを示した場面である（第 \ref{ch:event}・\ref{ch:text} 章）。
\item \textbf{読み取りをリスク管理に使う}。2 日の反応の\textbf{大きさ}を予測できるなら、それは売買の向きの予測でなくても、発表の前後の量の調整や、変動率の予測に価値がある（\chapref{ch:risk}）。
\item \textbf{コストと容量を明示する}。小型株では、平方根則（\chapref{ch:cost}）で、資金の規模ごとのコストを見積もり、「どの規模の資金で、この結論が成り立つか」を書く。
\item \textbf{先読みを避け続ける}。学習の期間の止まった言語モデル（ChronoGPT）や、会社名を伏せた入力（Glasserman・Lin 2023）で、LLM が将来を知っている影響を除く（第 \ref{ch:text} 章）。
\end{enumerate}

\section{してはいけないこと}

\begin{itemize}
\item \textbf{成績が出るまで条件を変えて試し続ける}。100 通りを試せば、無意味な戦略からでも年率 1 を超えるシャープレシオが見つかる（\chapref{ch:backtest}）。試した数を記録し、事前登録にない分析は探索として区別する。
\item \textbf{レバレッジやオプションの売りで、シャープレシオを良く見せる}。平時のシャープレシオは、まれな大損を示さない（第 \ref{ch:rv}・\ref{ch:failures} 章）。
\item \textbf{今日残っている会社だけで、過去の成績を評価する}。生存者バイアスは、年に数\% の大きさになりうる（第 \ref{ch:hfperf} 章）。LLM-BLM の小型株のデータでは、この偏りを別に見積もった。
\end{itemize}

\section{結び}

ヘッジファンドの研究から分かるのは、市場に勝つことは、プロにとっても非常に難しいということである。勝っている少数の運用者は、(1) 多くの弱いシグナルを組み合わせ、(2) 速く、安く売買し、(3) ほかの人が取りたがらないリスク（流動性、急落、混雑の解消）を、それに耐えられる資金の構造で引き受け、(4) 新しいデータと方法を探し続けている。一つの新しい道具（LLM）だけで市場の平均を上回ることは、文献から見て、期待しにくい。

一方で、研究としての価値は、勝てたかどうかだけにあるのではない。LLM が決算発表文の何を読み、その情報がどの速さで価格に入り、どこに遅れが残るのかを、先読みと過学習を避けて正しく測ることは、それ自体が市場の効率についての新しい証拠になる。
